\section{Disjoint witnesses and exclusion of corners}
\label{df:sec:witnesses}

We now turn the crossing estimates into an obstruction to a corner of the
limit shape. Assume that the two extreme supporting slopes in a fixed
direction differ. We will choose one deterministic scale $(s,L)$ at which
certain local path events are rare, and then divide every simple path from
the origin into chunks at that scale. A full chunk with small defect
supplies a witness for one of these events. More generally, the number of
disjoint witnesses in a chunk controls how negative its defect can be.

These deterministic bounds force any long path with small total defect to
contain many edge-disjoint witnesses. The BK inequality and a count of
coarse itineraries make the existence of such paths summably unlikely.
At the chosen scale this almost surely yields a positive average defect
per full chunk, contradicting the shape theorem along the direction of
contact. The argument allows the paths, the chunks, and their witnesses to
be chosen after the weights are observed.

Kesten's path-length estimates \cite[Section~3]{Kesten1980} already
count rectangle itineraries, record repeated rectangles through multiple
vertex-disjoint crossings, and use independence over spatially disjoint
rectangles. Here the witnesses certify negative excess relative to a
supporting functional, and BK controls their disjoint occurrence even
when the coarse test regions overlap.

\subsection{Disjoint occurrence for strict path witnesses}

The required form of BK concerns the edges that certify the events.
Test regions may overlap, and the same event may appear several times in
the list. This is essential when a path returns to a coarse cell.

Continuous-product extensions of disjoint occurrence are treated more
generally by Arratia, Garibaldi, and Hales \cite{AGH2018}. We give the
finite-bit reduction for the strict path events used here.

\begin{lemma}[BK for strict path witnesses]\label{df:wi:continuous-bk}
Let $A_1,\ldots,A_q$ be a finite list of events. For each $j$, suppose
$A_j$ is the existence of a path $P$ in a prescribed family of finite
lattice paths, all of whose edges belong to a fixed finite edge set,
satisfying
\[
 \tau(P)<b_j(P),
\]
where $b_j(P)$ is deterministic. A family may include several labeled
copies of a path with different bounds. Let $D(A_1,\ldots,A_q)$ be
the event that one can choose a witnessing path for each event with
pairwise disjoint edge sets. Then
\begin{equation}\label{df:wi:bk-bound}
 \PP(D(A_1,\ldots,A_q))\le\prod_{j=1}^q\PP(A_j).
\end{equation}
The events in the list need not be distinct.
\end{lemma}
\begin{proof}
Let $F_{\kappa,\lambda}$ be the Gamma distribution function and define
its quantile function by
\[
 \Phi(u)=F_{\kappa,\lambda}^{-1}(u)\quad(0\le u<1),
 \qquad\Phi(1)=+\infty.
\]
Here $\Phi(0)=0$. The density is positive for every $t>0$, so the
distribution function is continuous and strictly increasing from zero
to one. Consequently $\Phi$ is increasing and continuous on $[0,1)$.
Represent the independent weights as
\[
 \tau_e=\Phi(U_e),\qquad
 U_e=\sum_{h\ge1}2^{-h}B_{e,h},
\]
where all $B_{e,h}$ are independent Bernoulli variables of parameter
$1/2$. At precision $t$ put
\[
 U_e^{(t)}=2^{-t}+\sum_{h=1}^t2^{-h}B_{e,h},
 \qquad \tau_e^{(t)}=\Phi(U_e^{(t)}).
\]
Then $U_e^{(t)}\downarrow U_e$; because $U_e<1$ almost surely,
continuity of $\Phi$ gives $\tau_e^{(t)}\downarrow\tau_e$ almost surely
on the finite union of the edge sets under consideration. Define
$A_j^{(t)}$ and $D^{(t)}$ by using these rounded weights.

At each fixed precision, the events $A_j^{(t)}$ are decreasing events
of finitely many independent bits. Fixing all precision-$t$ bits of
the edges of a witnessing path forces its strict cost inequality
regardless of the other bits. Edge-disjoint path witnesses thus give
disjoint bit certificates. Complement the bits so that the events are
increasing. For each certificate, retain its realized one-bits and set
every other bit to zero. This configuration still belongs to the event,
because it agrees with the realized configuration on the certificate,
and its one-set lies inside that certificate. Disjoint certificates
therefore give disjoint positive witnesses in the convention of
\cite{BK85}; any unused one-bits can be assigned by monotonicity.
The finite Bernoulli BK inequality
\cite[Theorem~3.3, Equation~(3.4)]{BK85}, iterated over the list, gives
\[
 \PP(D^{(t)})\le\prod_{j=1}^q\PP(A_j^{(t)})
                   \le\prod_{j=1}^q\PP(A_j).
\]
Monotone disjoint occurrence is associative, so iteration allows
repeated events. The second inequality follows from
$\tau_e^{(t)}\ge\tau_e$.

The events $D^{(t)}$ increase with $t$. A fixed finite list of
original witnessing paths involves finitely many traversals. Its
rounded path costs decrease to the original costs, so every strict
inequality holds at all sufficiently fine precisions. Conversely,
every rounded witness remains a witness with the original weights.
Therefore, up to a null set,
\[
 D(A_1,\ldots,A_q)=\bigcup_{t\ge1}D^{(t)}.
\]
Continuity of probability from below proves \eqref{df:wi:bk-bound}.
No bound on the Gamma density or on the derivative of its quantile
function is used, so the proof includes $0<\kappa<1$.
\end{proof}

\subsection{Choosing one scale under a corner assumption}

Fix $b\in[-1,1]$ and use the notation of Lemma~\ref{df:prelim:supports}.
For the rest of the construction suppose, towards a contradiction, that
\[
  c=\frac{a_+-a_-}{2}>0,
  \qquad
  \eta=\min\{1,c/4\}.
\]
Choose $q_0\in(0,1)$ so small that
\begin{equation}
\label{df:wi:q-choice}
  400(2q_0)^{1/16}<\frac12.
\end{equation}
Use this value of $q_0$ in Lemmas~\ref{df:gr:characteristic-lengths}
and~\ref{df:gr:record-scales}, and write $s=2^k$ and $L=L_k$ along the resulting
subsequence.  In particular, $s$ and $L$ are integers and $L/s\to\infty$.

For $(i,h)\in\Z^2$, set
\[
  v_{i,h}=(iL,\lfloor biL\rfloor+hs),
  \qquad
  \mathcal B_{i,h}
  =v_{i,h}+\{z\in\Z^2:\abs{x(z)}\le4L,\ \abs{r(z)}\le4s\}.
\]
At each index $(i,h)$ we test for either a saving below an extreme support
or a cheap forward crossing. Precisely, define $\mathcal A_{i,h}$ to be
the union of the following events:
\begin{enumerate}
\item Some path $P$ with all its vertices in $\mathcal B_{i,h}$ satisfies
      $H_\sigma(P)<-\eta s$ for at least one sign
      $\sigma\in\{-,+\}$.
\item The event $E(s,L)$ holds at the lattice translate $v_{i+1,h}$.
\end{enumerate}
These are decreasing events of finitely many edge weights.  Throughout this
section a \emph{path witness} for $\mathcal A_{i,h}$ means a path satisfying
one of these two defining conditions, including its strict cost inequality.
All such events have the form in Lemma~\ref{df:wi:continuous-bk}.

The first event has probability $o(1)$ by
Lemma~\ref{df:gr:variable-endpoints}, and the second has probability at most
$q_0$ by the definition of $L$.  Translation invariance makes these bounds
uniform in $(i,h)$.  We may therefore fix one sufficiently large index $k$
on the subsequence such that
\begin{equation}
\label{df:wi:rare-events}
  \PP(\mathcal A_{i,h})\le2q_0
  \qquad\text{for every }(i,h)\in\Z^2.
\end{equation}
We also take this same index large enough that $s\ge2$ and
\begin{align}
\label{df:wi:deterministic-scale}
  C_f&:=\max\{\abs{m_0-ab},\abs{a}\}\le\eta s,\\
  2m_0L-\abs{a}(s+1)&>\eta s,
  \qquad
  3\eta s+2\abs{a}(s+1)<Ks.
  \nonumber
\end{align}
The first two requirements follow from $s\to\infty$ and $L/s\to\infty$;
the third follows from the definition of $K$ and $\eta\le1$.
The scale $k$, and hence $s$ and $L$, stays fixed from now on. The
probability bound \eqref{df:wi:rare-events} holds uniformly at this one
scale; only the length of the path will subsequently tend to infinity.
This is the finite-scale obstruction we will establish: a positive slope
gap and these rare local events almost surely give a positive lower bound
on $H(P)/N$ for every simple path from the origin with sufficiently many
full chunks $N$. In particular, $s/L$ is now a fixed positive number, and
the final contradiction uses this positivity.

\subsection{Extracting witnesses from path chunks}

Let $P$ be a finite simple lattice path starting at the origin.  Decompose
it into chunks as follows.  At the start of a chunk, with coordinates
$(x_0,r_0)$, follow the path until it first reaches a vertex for which
\[
  \abs{x-x_0}\ge2L
  \quad\text{or}\quad
  \abs{r-r_0}\ge s.
\]
Cut there and start the next chunk.  A chunk ending at such a cut is called
\emph{full}.  The decomposition consists of $N$ full chunks and a final
remainder, which may be empty.  Associate to each chunk, including the
remainder, the index
\begin{equation}
\label{df:wi:chunk-index}
  (i,h)=\bigl(\lfloor x_0/L\rfloor,\lfloor r_0/s\rfloor\bigr).
\end{equation}

The cuts and indices depend on the candidate path. All the estimates
below hold for each path individually and require no independence of its
chunks.

For a nearest-neighbor step, $\abs{\Delta x}\le1$ and
$\abs{\Delta r}\le1$.  Thus every vertex of a chunk has
$\abs{x-x_0}\le2L$ and $\abs{r-r_0}\le s+1$.  Moreover,
\[
  0\le x_0-iL<L,
  \qquad
  0\le r_0-r(v_{i,h})<s+1.
\]
The entire chunk, and each of its subpaths, therefore lies in
$\mathcal B_{i,h}$: its longitudinal offsets from $v_{i,h}$ belong to
$[-2L,3L)$, and its transverse offsets have absolute value at most
$2s+2\le4s$.

\begin{lemma}[A cheap full chunk supplies a witness]
\label{df:wi:chunk-witness}
If a full chunk $Q$ with index $(i,h)$ satisfies $H(Q)<\eta s$, then its
edges contain a path witness for $\mathcal A_{i,h}$.
\end{lemma}

\begin{proof}
There are three possible exit behaviors. A transverse exit uses the slope
gap, a backward exit has large defect by nonnegativity, and a forward exit
reduces to the crossing event after controlling its outer pieces.

If $\abs{r(\Delta Q)}\ge s$, the support identity gives
\[
  \min_{\sigma\in\{-,+\}}H_\sigma(Q)
  =H(Q)-c\abs{r(\Delta Q)}
  <(\eta-c)s\le-3\eta s<-\eta s.
\]
The whole chunk is then a witness for the first defining event.

Otherwise a longitudinal exit occurred, and
$x(\Delta Q)\in\{-2L,2L\}$.  In the negative case, nonnegative weights
and \eqref{df:wi:deterministic-scale} imply
\[
  H(Q)\ge2m_0L-\abs{a}(s+1)>\eta s,
\]
contrary to the hypothesis.

It remains to consider $x(\Delta Q)=2L$.  The initial column is strictly
to the left of $(i+1)L$, and the final column is at or to the right of
$(i+2)L$.  Let $z$ be the first vertex of $Q$ on column $(i+2)L$, and let
$p$ be the last vertex on column $(i+1)L$ preceding $z$.  The subpath
$Q'$ from $p$ to $z$ stays between these two columns.  Its transverse
offsets from $v_{i+1,h}$ have absolute value at most $2s+2\le8s$.
Consequently $Q'$ fits inside the translate of the block defining
$E(s,L)$ at $v_{i+1,h}$.

Write $Q=Q_{\rm in}*Q'*Q_{\rm out}$, allowing an empty outer part.
If either outer part has $H<-\eta s$, it witnesses the first defining
event, because it lies in $\mathcal B_{i,h}$ and
$\min_\sigma H_\sigma\le H$.  Otherwise additivity gives
$H(Q')<3\eta s$.  Since
$\abs{r(\Delta Q')}\le2(s+1)$, we obtain
\[
  \tau(Q')
  <m_0L+3\eta s+2\abs{a}(s+1)
  <m_0L+Ks.
\]
Thus $Q'$ witnesses the second defining event.
\end{proof}

\subsection{Controlling negative defect by witness multiplicity}

To control chunks whose defects are very negative, we count how many
edge-disjoint witnesses their edges support. A stopping rule will split
off a new witness each time a prefix accumulates a defect below $-\eta s$;
the bound on a single step limits the defect lost at each split.

For any chunk $Q$ with index $(i,h)$, let $M_Q$ be the maximum number of
pairwise edge-disjoint path witnesses for $\mathcal A_{i,h}$ whose edges
all belong to $Q$.  This maximum is finite.  Every witness uses at least
one edge: the first defining event has a strictly negative defect bound,
and the second crosses two distinct columns.  Hence $M_Q$ is bounded by
the number of edges of $Q$.

\begin{lemma}[Defect bound from disjoint witnesses]
\label{df:wi:chunk-defect}
Every chunk, including the remainder, satisfies
\begin{equation}
\label{df:wi:basic-defect}
  H(Q)\ge-\eta s-2\eta s M_Q.
\end{equation}
A full chunk satisfies the stronger bound
\begin{equation}
\label{df:wi:full-defect}
  H(Q)\ge\eta s-4\eta s M_Q.
\end{equation}
If $M_{\rm tot}$ denotes the sum of $M_Q$ over the $N$ full chunks and
the remainder of $P$, then
\begin{equation}
\label{df:wi:total-defect}
  H(P)\ge(N-1)\eta s-4\eta s M_{\rm tot}.
\end{equation}
The path contains $M_{\rm tot}$ pairwise edge-disjoint witnesses, assigned
to the corresponding chunk indices, even if some indices repeat.
\end{lemma}

\begin{proof}
For each oriented nearest-neighbor step $e$,
\[
  H(e)=\tau_e-f(\Delta e)\ge-C_f.
\]
Starting at the beginning of $Q$, split off the first prefix with
$H<-\eta s$, if one exists, and repeat with the remaining path.
At the last step of each split prefix, its defect decreases from a value
at least $-\eta s$ to a value below $-\eta s$.  By
\eqref{df:wi:deterministic-scale}, the defect of this prefix is at least
$-\eta s-C_f\ge-2\eta s$.  Every split prefix is a witness for the
original chunk's event $\mathcal A_{i,h}$: it remains in the same box,
and $\min_\sigma H_\sigma\le H<-\eta s$.

If this procedure splits off $t$ prefixes, simplicity of $P$ makes their
edge sets disjoint, so $t\le M_Q$.  The unsplit final part has
$H\ge-\eta s$.  Additivity now gives
\[
  H(Q)\ge-\eta s-2\eta s t
       \ge-\eta s-2\eta s M_Q,
\]
which proves \eqref{df:wi:basic-defect}.

For a full chunk with $M_Q=0$, Lemma~\ref{df:wi:chunk-witness} gives
$H(Q)\ge\eta s$.  If $M_Q\ge1$, then
\[
  -\eta s-2\eta s M_Q
  \ge\eta s-4\eta s M_Q.
\]
This proves \eqref{df:wi:full-defect}.  Apply that bound to all full chunks
and use \eqref{df:wi:basic-defect}, weakened to
$-\eta s-4\eta sM_Q$, on the remainder.  Their sum is
\eqref{df:wi:total-defect}.

Finally, choose a family attaining $M_Q$ in each chunk.  Distinct chunks
of a simple path have disjoint edge sets, so the union of these families
is still edge-disjoint.  The assertion does not require their associated
events or their boxes to be distinct.
\end{proof}

\subsection{Counting itineraries and witness multiplicities}

We now apply BK to the witnesses supplied by the deterministic bounds.
The counting records the coarse itinerary and the number of witnesses
assigned to each of its positions. For each fixed number of chunks this
produces a finite collection of deterministic event lists covering all
possible choices of the path and its witnesses, including paths that
revisit a coarse cell.

\begin{proposition}[Uniform lower bound for long paths]
\label{df:wi:long-paths}
At the fixed scale chosen above, almost surely there exists a finite
$N_0$ such that every finite simple path $P$ starting at the origin and
having $N\ge N_0$ full chunks satisfies
\begin{equation}
\label{df:wi:long-path-bound}
  H(P)\ge\frac12 N\eta s.
\end{equation}
\end{proposition}

\begin{proof}
First quantify how many witnesses a path of small defect must supply.
Let $B_N$ be the event that there exists such a simple path with $N$ full
chunks and $H(P)<N\eta s/2$.  If it occurs, \eqref{df:wi:total-defect}
implies
\[
  M_{\rm tot}>\frac N8-\frac14.
\]
For all sufficiently large $N$, we may therefore select exactly
$t_N=\lfloor N/16\rfloor$ of the edge-disjoint witnesses supplied by
Lemma~\ref{df:wi:chunk-defect}.

Next count the possible event lists. List the chunk indices in order,
including the index of the remainder.
The first is $(0,0)$.  Each coordinate of consecutive indices differs
by at most $3$, by the displacement bounds for a full chunk and
\eqref{df:wi:chunk-index}.  There are consequently at most $49^N$ possible
index itineraries.  For a fixed itinerary, the numbers of selected
witnesses in its $N+1$ positions form a weak composition of $t_N$.
The number of these compositions is
\[
  \binom{N+t_N}{t_N}\le2^{N+t_N}\le4^N.
\]
The allocation is indexed by positions along the itinerary, so repeated
visits to one coarse cell retain separate positions. Each composition
specifies a list of $t_N$ events
$\mathcal A_{i,h}$, with repetitions when appropriate.  Its selected
witnesses certify disjoint occurrence of that list.  Lemma~\ref{df:wi:continuous-bk}
and \eqref{df:wi:rare-events} bound its probability by $(2q_0)^{t_N}$.
There is no additional count of witness paths: that lemma bounds the
existence of all choices of disjoint witnessing paths at once.
Thus the path-dependent selection is covered by a union over deterministic
itineraries and compositions. Taking this union yields
\begin{equation}
\label{df:wi:entropy-bound}
  \PP(B_N)\le196^N(2q_0)^{\lfloor N/16\rfloor}
\end{equation}
for all sufficiently large $N$.  All the events involved are measurable;
in particular, finite lattice paths form a countable family.

Put $p=2q_0$.  Condition~\eqref{df:wi:q-choice} implies $0<p<1$ and
$196p^{1/16}<1$, while the right-hand side of
\eqref{df:wi:entropy-bound} is at most
\[
  p^{-1}\bigl(196p^{1/16}\bigr)^N.
\]
Thus $\sum_N\PP(B_N)<\infty$.  The first Borel--Cantelli lemma proves
the simultaneous eventual bound \eqref{df:wi:long-path-bound} for all the
paths in the proposition.
\end{proof}

\subsection{Contradiction with the shape theorem}

The bound just proved applies simultaneously to all sufficiently long
simple paths from the origin. In particular it applies to approximate
minimizers selected from the realized environment. Along the contact
direction $u$, their number of chunks must grow with their longitudinal
displacement, whereas the shape theorem makes their defect sublinear in
that displacement.

\begin{proposition}[Absence of corners]
\label{df:wi:no-corner}
For every $b\in[-1,1]$, the two one-sided transverse derivatives at
$u=(1,b)$ agree:
\[
  a_-=a_+.
\]
\end{proposition}

\begin{proof}
Fix an arbitrary deterministic $b\in[-1,1]$ and suppose that $c>0$.
Make all the preceding choices for this $b$, including the single fixed
scale $k$.  For positive integers $t$, let
\[
  z_t=(t,\lfloor bt\rfloor).
\]
A path from $0$ to $z_t$ with $N$ full chunks has $N+1$ chunks when the
remainder is included.  Each changes the $x$-coordinate by at most $2L$,
so
\begin{equation}
\label{df:wi:cut-number}
  N+1\ge\frac{t}{2L}.
\end{equation}

Work on the probability-one event where the shape theorem and
Proposition~\ref{df:wi:long-paths} both hold.  For each $t$, choose a finite
path from $0$ to $z_t$ with cost at most $T(0,z_t)+1$ and erase its loops.
The resulting simple path $P_t$ still has cost between $T(0,z_t)$ and
$T(0,z_t)+1$, since the edge weights are nonnegative.  The shape theorem
gives $T(0,z_t)/t\to m_0$, and
\[
  f(z_t)=m_0t+a(\lfloor bt\rfloor-bt)=m_0t+O(1).
\]
It follows that
\begin{equation}
\label{df:wi:shape-defect-zero}
  \frac{H(P_t)}{t}\longrightarrow0.
\end{equation}

On the other hand, the numbers $N_t$ of full chunks of $P_t$ tend to
infinity by \eqref{df:wi:cut-number}.  Proposition~\ref{df:wi:long-paths}
therefore applies for all sufficiently large $t$ and gives
\[
  \liminf_{t\to\infty}\frac{H(P_t)}t
  \ge\frac{\eta s}{2}
       \liminf_{t\to\infty}\frac{N_t}{t}
  \ge\frac{\eta s}{4L}>0.
\]
This contradicts \eqref{df:wi:shape-defect-zero}.  Hence $c>0$ is impossible,
and convexity gives $a_-=a_+$.

The choice of $b$ was arbitrary and deterministic.  The corner gap of the
deterministic norm is itself deterministic, and the argument excludes a
positive gap for each such $b$.  Accordingly the conclusion holds for
every $b\in[-1,1]$; no intersection of uncountably many probability-one
events is being asserted or needed.
\end{proof}
